\begin{tcolorbox}[colback=white, colframe=my-green, coltitle=black, parskip=5mm, title=\textbf{Next-Hop Question, Fixed Target (Short)}, halign=flush left]
\begin{scriptsize}
\ttfamily
\setlength{\parindent}{0pt}
\normalfont\textit{Short paper variant; in-context examples omitted for space.}\par
\medskip
\ttfamily
Write the next question in this multi-hop QA chain.\par
\mbox{}\par
Previous question:\par
\textcolor{my-green}{\textbf{\{prev\_question\}}}\par
\mbox{}\par
Previous answer:\par
"\textcolor{my-green}{\textbf{\{prev\_answer\}}}"\par
\mbox{}\par
Target answer:\par
"\textcolor{my-green}{\textbf{\{target\_answer\}}}"\par
\mbox{}\par
Source excerpt:\par
\textcolor{my-green}{\textbf{\{source\_excerpt\}}}\par
\mbox{}\par
Current document:\par
\textcolor{my-green}{\textbf{\{doc\_text\}}}\par
\mbox{}\par
Write one simple question that uses "\textcolor{my-green}{\textbf{\{prev\_answer\}}}" exactly and asks for one new fact from this document.\par
Think through a few possible questions first, then choose the best honest one.\par
Do not force a relationship. If the document only co-mentions the values, or if the hop would require guesswork, arithmetic, or a made-up connection, output NO\_VALID\_QUESTION.\par
Avoid generic "what year/date/quarter is associated with..." questions unless that date is truly the fact being asked for.\par
If the answer is a date or time, use the full supported granularity rather than a fragment.\par
The question must uniquely identify the answer from the quoted evidence. If multiple people, places, sections, dates, or values in the excerpt could fit, output NO\_VALID\_QUESTION.\par
Write exactly one sentence total. Do not write a setup sentence followed by a second question or any extra sentence fragment.\par
Use the previous answer as an inline part of that single sentence, not as a separate lead-in.\par
The goal is to use the previous answer to identify the right subject, event, or record in this document, then ask for exactly one new fact about that identified subject.\par
If the answer is a named entity, ask for the full entity, not a partial fragment like a surname, section numeral, or single generic word unless the source itself only gives that fragment.\par
Prefer distinctive, interesting answers and questions over low-information ones. Avoid generic targets like a bare year, quarter, section label, ordinal word, or generic capitalized word unless that is clearly the central fact.\par
Make the question hard to misread as being about some other person, event, or document that also mentions the same bridge.\par
Use only the minimum detail needed to uniquely identify the subject. Do not pile on extra descriptors once the subject is already pinned down.\par
As a soft preference, favor entities and facts that are broadly meaningful or likely to matter to a company researcher: organizations, products, policies, markets, hospitals, countries, technologies, roles, or durable metrics.\par
Avoid dead-end trivia when possible, especially random individual people, one-off titles, or isolated labels that are unlikely to support a useful next hop, unless the person or label is clearly the main subject of the excerpt.\par
\mbox{}\par
\textcolor{my-green}{\textbf{\textless{}\textless{}ICL\_EXAMPLES\_OMITTED\_FOR\_SPACE\textgreater{}\textgreater{}}}\par
\mbox{}\par
End with exactly one final block using these labels:\par
QUOTE: \textless{}exact supporting excerpt(s) from the document; use " | " if needed\textgreater{}\par
REASONING: \textless{}brief explanation of why this is the best honest option; keep this short and do not reveal long chain-of-thought\textgreater{}\par
QUESTION: \textless{}one single-line question containing "\textcolor{my-green}{\textbf{\{prev\_answer\}}}"\textgreater{}\par
ANSWER: \textcolor{my-green}{\textbf{\{target\_answer\}}}\par
\end{scriptsize}
\end{tcolorbox}
